%%%PAGE 289 rot=0 legibility=good summary=斜面上圆周运动中重力功率、两球关联速度、质点位移相同的时间间隔、磁场中带电粒子的洛伦兹力分解（动量定理与动能定理）、环套杆摩擦做功三种情形、不同倾角斜面下滑时间相等（页面左缘有装订/贴纸痕迹）
%%%BLOCK id=289-01 ch=04 sec=斜面上的圆周运动·重力功率 kind=note alt=06 cont=none y=0.00-0.13
%FIG p289_a 0.095 0.0 0.43 0.125
\notefig[0.3]{figures/p289_a.png}

$P_G$ 先↑后↓

$P_{G\max}$ 时 $F_{\text{合}}$ 与 $OA$ 平行
%%%END
%%%BLOCK id=289-02 ch=04 sec=运动的合成与分解·关联速度 kind=problem alt=none cont=none y=0.13-0.27
%FIG p289_b 0.105 0.135 0.29 0.26
\notefig[0.3]{figures/p289_b.png}

%FIG p289_c 0.325 0.115 0.65 0.258
\notefig[0.4]{figures/p289_c.png}

$V_A\sin\theta=V_B\cos\dfrac{\theta}{2}$

两球速度大小相等时 $B$ 在 $O$ 点正上方% CHECK: “O点”手写像 O（图中对应点标注像 D/O）
%%%END
%%%BLOCK id=289-03 ch=08 sec=机械振动·质点位移相同 kind=note alt=none cont=none y=0.28-0.37
\unclear{无搅3算}\quad $x^2=g$% CHECK: 首行第2、4字不清（第4字被反复涂写）；“x”也可能是 $\omega$

每隔 $t$ 两\unclear{个}质点位移相同一次\quad $2t=nT$
%%%END
%%%BLOCK id=289-04 ch=11 sec=洛伦兹力分解·动量定理 kind=problem alt=06 cont=none y=0.38-0.60
%FIG p289_d 0.055 0.38 0.25 0.495
\notefig[0.28]{figures/p289_d.png}

\begin{gather*}
-mgt-qBV_xt=0-mV_0\\
mgt+\frac12 Bql=\text{\struck{$0-$}}\,mV_0\qquad t=\frac{mV_0-\frac12 qBL}{mg}\\
-mgd=\frac12 m\left(V_{\text{末}}^2-V_0^2\right)
\end{gather*}
% CHECK: 第三式 V 的下标手写像“末”；第二式中 l 与 t 式中 L 应为同一量（图中标 L/2）；第二式等号后被划去的是“0-”
%%%END
%%%BLOCK id=289-05 ch=11 sec=洛伦兹力·环套杆摩擦做功 kind=problem alt=06 cont=none y=0.61-0.78
%FIG p289_e 0.10 0.635 0.27 0.72
\notefig[0.28]{figures/p289_e.png}

%FIG p289_f 0.295 0.635 0.86 0.78
\notefig[0.55]{figures/p289_f.png}

\begin{itemize}
\item $mg=qVB$\quad $W_f=0$
\item $mg\le qVB$\quad $W_f=\text{\unclear{…}}-\dfrac12\left(mV_0^2-V_{\text{末}}^2\right)$
\item $mg>qVB$\quad $W_f=\dfrac12 mV_0^2$
\end{itemize}
% CHECK: 第二种情形 W_f 等号后有一个无法辨认的涂改小记号（像被划去的 E/½）；括号内第二项漏写 m；V 的下标手写像“末”；“mg≤qVB”也可能是 mg<qVB
%%%END
%%%BLOCK id=289-06 ch=03 sec=斜面·不同倾角等时下滑 kind=problem alt=none cont=none y=0.78-0.99
%FIG p289_g 0.115 0.785 0.36 0.89
\notefig[0.3]{figures/p289_g.png}

\[ t=\sqrt{\frac{2L}{\text{\struck{$\sin\theta$}}\cos\theta\, g\sin\theta}}=\sqrt{\frac{2L}{\frac12\sin2\theta\, g}} \]
% CHECK: 第一个根号分母中被划去的字形不清（像 sinθ）
\[ t_1=t_2 \]
%%%END
%%%PAGEEND 289
